\documentclass[11pt,a4paper]{article}
\usepackage[T1]{fontenc}
\usepackage[utf8]{inputenc}
\usepackage{lmodern}
\usepackage[margin=24mm,headheight=15pt]{geometry}
\usepackage{amsmath,amssymb,amsthm,mathtools}
\usepackage{booktabs,longtable,array}
\usepackage[dvipsnames]{xcolor}
\usepackage{microtype}
\usepackage{enumitem}
\usepackage{fancyhdr}
\usepackage{hyperref}
\usepackage{xurl}
\hypersetup{colorlinks=true,linkcolor=MidnightBlue,citecolor=MidnightBlue,urlcolor=MidnightBlue,
  pdftitle={Explicit Diophantine formulations of Con(ZFC)},
  pdfauthor={Research contribution prepared for Vladimir Reshetnikov / ProveIt},
  pdfsubject={K215 specialization, fixed-marker reduction, explicit polynomial systems}}
\setlength{\parskip}{5pt}
\setlength{\parindent}{0pt}
\setcounter{tocdepth}{1}
\setlist{nosep,leftmargin=1.5em}
\pagestyle{fancy}
\fancyhf{}
\fancyhead[L]{\small Explicit Diophantine formulations of $\operatorname{Con}(\mathrm{ZFC})$}
\fancyhead[R]{\small ProveIt research contribution}
\fancyfoot[C]{\thepage}
\newtheorem{theorem}{Theorem}[section]
\newtheorem{lemma}[theorem]{Lemma}
\newtheorem{proposition}[theorem]{Proposition}
\theoremstyle{remark}
\newtheorem{remark}[theorem]{Remark}
\newcommand{\Z}{\mathbb Z}
\newcommand{\N}{\mathbb N}
\newcommand{\Zp}{\mathbb Z_{>0}}
\newcommand{\Con}{\operatorname{Con}}
\newcommand{\ZFC}{\mathrm{ZFC}}
\newcommand{\ZRef}{\operatorname{Ref}_{\ZFC}}
\newcommand{\code}[1]{\nolinkurl{#1}}
\newcommand{\pathcode}[1]{\texttt{\small\detokenize{#1}}}
\newcommand{\repo}{\url{https://github.com/VladimirReshetnikov/ProveIt}}
\newcommand{\commit}{\code{cdd582bec13f908df3711fa9e05d1ddac368a5db}}
\newcommand{\mc}{M_C}
\newcommand{\mf}{M_F}
\newcommand{\nf}{N_1}
\newcommand{\nm}{N_2}
\newcommand{\nw}{N_W}
\newcommand{\nk}{N_k}
\newcommand{\nt}{N_t}
\allowdisplaybreaks[2]

\begin{document}
\hypersetup{pageanchor=false}
\begin{titlepage}
\vspace*{12mm}
{\color{MidnightBlue}\Large PROVEIT\quad /\quad RESEARCH CONTRIBUTION\par}
\vspace{14mm}
{\Huge\bfseries Explicit Diophantine\par formulations of\par $\operatorname{Con}(\mathrm{ZFC})$\par}
\vspace{9mm}
{\Large A fixed K215 specialization, a seven-equation system,\par
and a polynomial in fifteen positive integer variables\par}
\vspace{12mm}
{\large Prepared for Vladimir Reshetnikov\par 22 October 2026\par}
\vspace{12mm}
\begin{minipage}{0.96\textwidth}
\textbf{Recommended readable form.} Seven polynomial equations in sixteen
positive integer variables, maximum degree $28$, with a shared $64$-operation
evaluation circuit. Their sum of squared residuals is one polynomial of
degree $56$, evaluated in $84$ operations.

\textbf{Fewer variables.} One polynomial in fifteen positive integer
variables, degree $295$, evaluated in $82$ operations.

\textbf{Meaning.} Each formulation has a solution exactly when the fixed
recursive proof system for ZFC has a finite refutation. Consequently its
lack of solutions is equivalent, by mathematics formalizable in ZFC, to
$\Con(\ZFC)$.

\textbf{Scope.} This is an explicit mathematical construction with
reproducible source and algebra audits. The inherited K215 machine endpoint
has Lean and Rocq verification. The new Diophantine reduction and the
inherited fixed-program polynomial compiler are not newly kernel checked
by this submission. No claim of global minimality is made.
\end{minipage}
\vfill
\small Source snapshot: \commit.\par
All coefficients are fixed by finite formulas and the included literal
machine table. Their enormous decimal expansions are not materialized.
\end{titlepage}
\hypersetup{pageanchor=true}

\begin{abstract}
We specialize the ProveIt fixed-program Diophantine compiler to the exact
215-state K215 machine calibrated to ZFC refutation. Since only one input
is needed, the universal construction's input Pell block can be replaced
by equality to a fixed marker. We prove that the three deleted input
witnesses can be restored as positive integers after main-rank recovery;
the proof uses neither a decoded computation nor the power-recovery
lemma. Combining this reduction with the repository's ordinary auxiliary
projection gives seven equations in sixteen positive unknowns. Combining
it with the normalized strong-root elimination gives one polynomial in
fifteen positive unknowns. We specify an actual coefficient slice by a
finite machine-to-window compiler, supply exact evaluation circuits,
prove uniform degree bounds, and give unrestricted-integer and quartic
variants. The distinction between a mathematical proof in ZFC, an
external proof-assistant theorem, and a new end-to-end formalization is
kept explicit throughout.
\end{abstract}

\tableofcontents
\newpage

\section{Result, domain, and meaning of simplicity}
\label{sec:result}

Fix the recursively presented ZFC proof calculus used by the repository's
K215 endpoint. Let $\ZRef(p)$ mean that $p$ codes a finite derivation of
contradiction in this calculus, and put
\[
\Con(\ZFC)\;:\!\iff\;\neg\exists p\in\N\,\ZRef(p).
\]
The source endpoint is the theorem
\code{KEndpoint215.blank_halts_iff_zfc_refutation}; its refutation
proposition is \code{BProv ZFCax_s [] Form.fBot} \cite{repo,k215}.
Coding finite syntax and finite derivations by natural numbers identifies
this proposition with $\exists p\,\ZRef(p)$.

\begin{theorem}[Explicit fixed-instance result]
\label{thm:main}
With the fixed integers defined in Section~\ref{sec:constants}, the system
$\mathcal S_{16}$ in Section~\ref{sec:system} and the polynomial $P_{15}$
in Section~\ref{sec:p15} satisfy
\begin{align*}
\ZFC&\vdash
 \bigl(\exists \mathbf z\in\Zp^{16}\;\mathcal S_{16}(\mathbf z)\bigr)
 \;\longleftrightarrow\;\exists p\in\N\,\ZRef(p),\\
\ZFC&\vdash
 \bigl(\exists \mathbf z\in\Zp^{15}\;P_{15}(\mathbf z)=0\bigr)
 \;\longleftrightarrow\;\exists p\in\N\,\ZRef(p).
\end{align*}
Negating either equivalence gives the requested equivalence between lack
of solutions and $\Con(\ZFC)$. The same holds for the sum-of-squares and
domain-conversion variants below.
\end{theorem}

The theorem is an ordinary mathematical theorem, using the explicitly
identified inherited compiler and Pell lemmas. Its proof does not use an
assumption that ZFC is consistent. Section~\ref{sec:formal} explains the
formalization boundary; the included executable checks are not a
replacement for the mathematical proof.

\begin{center}
\begin{tabular}{lrrrr}
\toprule
Form & Positive unknowns & Equations & Degree & Operations\\
\midrule
Recommended $\mathcal S_{16}$ & 16 & 7 & 28 & 64\\
Sum of squared residuals & 16 & 1 & 56 & 84\\
Normalized product & 16 & 1 & 157 & 74\\
Reduced auxiliary product & 16 & 1 & 125 & 75\\
Ordinary auxiliary product & 16 & 1 & 101 & 76\\
$P_{15}$, eliminated strong root & 15 & 1 & 295 & 82\\
Eliminated-root comparison form & 15 & 6 & 244 & 71\\
\bottomrule
\end{tabular}
\end{center}

For a system, ``degree'' means the maximum degree of its residual
polynomials after all abbreviations are expanded. ``Operations'' means
binary additions, subtractions, and multiplications in the supplied
shared evaluation circuit, including multiplication by a nontrivial
constant. For a comparison system it counts computation of the comparison
sides. It does not count expansion or construction of fixed coefficients.
The circuits contain no division or variable exponentiation.

There is no single mathematical ordering called ``simplest.'' Number of
unknowns, degree, arithmetic circuit size, coefficient bit length, and
proof complexity compete. This submission favors a short displayed
system and a small fully specified circuit. It does not claim a minimum
over all possible Diophantine representations. In particular, classical
universal equations with nine nonnegative unknowns exist
\cite{jones1982}; the present fifteen-variable result is not a
variable-count record for general existence theorems.

\section{Fixed constants and shared polynomial abbreviations}
\label{sec:notation}

The fixed radix parameters are
\begin{equation}
 b=5^{42},\qquad L_c=5^{43},\qquad d=bL_c=5^{85},\qquad
 B=2^d,\qquad W_*=2^bB^2=2^{2d+b}.
\label{eq:radices}
\end{equation}
The other fixed integers $K,\mc,\mf$ are the finite sums in
Section~\ref{sec:constants}. In particular,
\begin{equation}
 K=D_C+BD_R,\qquad \mf=M_F^{\mathrm{native}}+B-1.
\label{eq:ports}
\end{equation}
The shift by $B-1$ in $\mf$ is essential. The source port called
\code{MF} is not the native field mask.

The following definitions are shared by the main constructions:
\begin{align}
 q&=(B-1)J+1,& X&=wq,& Y&=sq^3,& E&=XY,\label{eq:common1}\\
 k&=\eta+\zeta,& c&=kY+\eta,& a&=E+Y,& H&=4a+3,\label{eq:common2}\\
 \Delta&=a^2+H=(a+2)^2-1,&
 D_0&=X+ac+\gamma H,\label{eq:common3}\\
 C&=q-F-Z-\alpha-2d,& W&=C-Z,\label{eq:common4}\\
 R&=(q^2-Z-qF)(q^2-1)+(\mc+q\mf)J,& L_0&=E(kY).
 \label{eq:common5}
\end{align}
Define four useful factors
\begin{align}
 \nf&=\tau^2-L_0(L_0+k),&
 \nm&=D_0^2-\Delta c^2,\label{eq:factors1}\\
 \nk&=k-hE-R,&
 \nt&=(K+w)C+q-F-t(q-1).\label{eq:factors2}
\end{align}

Every letter on the left of these definitions is a polynomial
abbreviation, not an additional unknown. Only the explicitly listed
witnesses are quantified. Although those witnesses are positive, some
computed abbreviations, such as $C,W,R$, can be negative away from a
solution. Soundness proofs must establish the required bounds at a
solution rather than assume them from notation.

\section{The recommended seven-equation system}
\label{sec:system}

Take the sixteen strictly positive integer unknowns
\begin{equation}
 (J,F,\alpha,t,f,h,i,T,s,w,\tau,\eta,\zeta,y,Z,\gamma)\in\Zp^{16}.
\label{eq:16vars}
\end{equation}
Besides the shared abbreviations, put
\begin{equation}
 V=c(Tf-1)-Rf^2,\qquad K_a=\Delta(f^2-1).
\label{eq:ordinaryaux}
\end{equation}
Our recommended system is
\begin{equation}
\boxed{\quad
\begin{aligned}
 \tau^2-L_0(L_0+k)&=1,\\
 D_0^2-\Delta c^2&=1,\\
 (ic^2)^2&=K_a,\\
 K_a(V^2-y^2)+y^2&=1,\\
 k-hE-R&=1,\\
 (K+w)C+q-F-t(q-1)&=1,\\
 W&=W_*.
\end{aligned}\quad}
\label{eq:S16}
\end{equation}

This is literally a system over $\Z$ with sixteen positive unknowns; it
contains no unbounded quantifier disguised as an abbreviation. The
constants in its coefficients are fixed before any unknown is chosen.
The shared circuit has $34$ multiplications and $30$
additions/subtractions. Its maximum residual degree is $28$.

Let $G_1,\ldots,G_7$ be the left sides minus the right sides of
\eqref{eq:S16}, in displayed order. Then
\begin{equation}
 \boxed{\quad H_{16}=G_1^2+G_2^2+\cdots+G_7^2=0.\quad}
\label{eq:sos}
\end{equation}
Over integers this is equivalent to the seven equations. It has exact
degree $56$, and the supplied shared circuit uses $41$ multiplications
and $43$ additions/subtractions, for $84$ operations in total.

The witness $i$ here belongs to the ordinary strong equation. The
normalized construction used next has a different normalization of its
witness also called $i$. We assert equality of the accepted proposition
between these constructions, not coordinatewise equality of their
solution sets.

\section{A single polynomial in fifteen positive unknowns}
\label{sec:p15}

Take the fifteen strictly positive unknowns
\[
 (J,F,\alpha,t,h,i,T,s,w,\tau,\eta,\zeta,y,Z,\gamma)\in\Zp^{15}.
\]
Use \eqref{eq:common1}--\eqref{eq:factors2}, and define
\begin{align}
 \nw&=1+(W-W_*)^2,\label{eq:markerfactor}\\
 S&=\Delta i c^2,& Q&=S^2,& D_*&=1+\Delta i^2c^4,\label{eq:elim1}\\
 b_*&=c+RD_*,& u_*&=cT,\label{eq:elim2}\\
 A_*&=Q(D_*u_*^2+b_*^2)-(Q-1)y^2-1,&
 B_*&=2Qu_*b_*,\label{eq:elim3}\\
 N&=A_*^2-D_*B_*^2.\label{eq:elim4}
\end{align}
The requested single polynomial can be taken to be
\begin{equation}
 \boxed{\quad P_{15}=\nf\nm\nw\nk\nt\,(N+1)-1=0.\quad}
\label{eq:P15}
\end{equation}
Its degree is exactly $295$. The supplied circuit evaluates it in
$46$ multiplications and $36$ additions/subtractions. In particular,
the apparently complicated local norm is computed using
$D_*B_*^2=4(QD_*u_*^2)(Qb_*^2)$, sharing subexpressions already needed
for $A_*$.

A six-equation alternative imposes
\[
 \nf=\nm=\nk=\nt=1,\qquad W=W_*,\qquad A_*^2=D_*B_*^2.
\]
It has the same positive solvability, maximum degree $244$, and a
$71$-operation shared comparison circuit. For exposition the
seven-equation ordinary system is substantially easier to inspect.

\section{Correctness: the fixed-marker reduction}
\label{sec:proof}

\subsection{The inherited interface, stated explicitly}

The universal-input construction in the pinned repository represents
acceptance of a fixed normalized machine using six fixed integer ports
and positive witnesses. We use the complete75 compiler and its
complete84/85/86/92 projections \cite{compiler,pell,projections}.
Its ordinary input is a positive integer $x$; for its input convention
the unary word has length $2x+2$. Our normalized machine ignores that
word, so the choice $x=1$ is sufficient.

In the complete84 source the main Pell root has the form
\[
 D_0=X+ac+(\rho+\sigma)H,
\]
and the input block enforces, at a zero,
\[
 W=2^{2dx+b}.
\]
It uses three positive witnesses $\delta,\rho,\sigma$. The fixed-marker
operation replaces $\rho+\sigma$ by a free positive $\gamma$, deletes
the three witnesses, fixes $x=1$, and imposes $W=W_*$; in a product
polynomial the last equality is represented by $\nw=1+(W-W_*)^2$.
Two quantified witnesses are saved overall.

The key inherited lemma is local to the Pell kernel. Write
\[
 (A+\sqrt{A^2-1})^n=\chi_A(n)+\psi_A(n)\sqrt{A^2-1},\qquad A=a+2.
\]

\begin{lemma}[Input-independent rank recovery]
\label{lem:rank}
On the valid compiler slice of Section~\ref{sec:constants}, the retained
first and main Pell conditions, index and transport conditions, and
ordinary auxiliary conditions of $\mathcal S_{16}$ give
\begin{equation}
 R\text{ odd},\qquad D_0=\chi_A(R),\qquad c=\psi_A(R),
 \qquad R>q.
\label{eq:rank}
\end{equation}
The normalized auxiliary conditions
\begin{equation}
 f^2-\Delta i^2c^4=1,\qquad
 QV^2-(Q-1)y^2=1,\qquad Q=\Delta^2i^2c^4
\label{eq:normalized}
\end{equation}
give the same conclusion when $f>0$ and $T$ is any integer. In the
product versions, they also force $\nk=\nt=1$ after the initial
integer-unit argument.
\end{lemma}

This is the precise imported rank lemma, not an assertion that a finite
test establishes all Pell solutions. The proofs are supplied as the
pinned notes \code{review_complete85_auxiliary_bezout_math.md},
\code{complete84_signed_quotient_absorption.md}, and
\code{review_complete86_ordinary_auxiliary_math.md}; the independent
review in this archive traces which hypotheses each step uses.

For orientation, the raw packing and transport bounds established before
rank recovery are
\begin{equation}
 q\ge B\ge16,\qquad 3q+1<R,\qquad R+2<q^4\le E<a.
\label{eq:rawbounds}
\end{equation}
The first equation can be written
$\tau^2-U(U+1)k^2=1$, where $U=XY^2$, and yields
$k=2\psi_{2U+1}(n)$. The main equation gives
$c=\psi_A(p)$. Strong divisibility and the auxiliary congruences
then give $R\equiv\pm p\pmod c$; the proved range bounds force
$p=R$. The ratio $kY<c<k(Y+1)$ excludes the wrong index sign.
These arguments use the old input factor only as an integer unit with
value $+1$ and use $\rho+\sigma$ only through positivity. The marker
factor is positive and the new $\gamma$ is positive, so those hypotheses
are preserved.

For the ordinary strong equation, the proof specifically establishes
$p\mid m$ and $c\mid m$, where
$f=\chi_A(m)$ and $ic^2=\Delta\psi_A(m)$. It must not assume the
stronger normalized divisibility $c^2\mid\psi_A(m)$. The supplied
ordinary proof uses only the weaker, valid relations. With seven
separate equations the required factor values are imposed directly.

\subsection{Restoration of the deleted positive witnesses}

\begin{lemma}[Fixed-marker restoration]
\label{lem:restore}
Suppose the retained conditions give \eqref{eq:rank}, and $W=W_*$.
Then there exist positive integers $\delta,\rho,\sigma$ that restore
the literal original input block at $x=1$, with
$\rho+\sigma=\gamma$.
\end{lemma}

\begin{proof}
Put $u=2d+b$. This is a fixed odd integer at least $3$.
The marker and the definitions of $C,W$ give
\[
 q=F+2Z+\alpha+2d+W_*>W_*=2^u>u.
\]
Thus $R>u$, and since $R,u$ are odd, $R\ge u+2$. Define
\begin{equation}
 \kappa=\psi_A(u),\quad \mu=\chi_A(u),\qquad
 \delta=\frac{\kappa-u}{\Delta},\quad
 \rho=\frac{\mu-a\kappa-2^u}{H},\quad
 \sigma=\gamma-\rho.
\label{eq:restore}
\end{equation}
These divisions occur only in the proof of witness existence, not in
the polynomial.

Pell expansion modulo $\Delta=A^2-1$ gives
$\psi_A(j)\equiv jA^{j-1}\pmod\Delta$. Since $u$ is odd,
$\kappa\equiv u\pmod\Delta$. Strict growth gives
$\psi_A(u)>u$, so $\delta$ is a positive integer.

Let $E_A(j)=\chi_A(j)-a\psi_A(j)$. Its initial values are $1,2$,
and it satisfies
\[
 E_A(j+2)=2AE_A(j+1)-E_A(j).
\]
The sequence $2^j$ satisfies the same recurrence modulo $H=4a+3$,
because $4-4A+1=-H$. Hence $E_A(j)\equiv2^j\pmod H$.
Also $E_A(2)=4A-1>4$ and $E_A(j+1)>2E_A(j)$ for $j\ge1$,
so $E_A(u)>2^u$. Therefore $\rho$ is a positive integer.

It remains to prove $\sigma>0$; this is the substantive restoration
condition. From the main root,
$E_A(R)=X+\gamma H$, while $E_A(u)=2^u+\rho H$. Therefore
\begin{equation}
 H(\gamma-\rho)=E_A(R)-E_A(u)-X+2^u.
\label{eq:gammadiff}
\end{equation}
For $j\ge1$,
$E_A(j)=2\psi_A(j)-\psi_A(j-1)$, so
$\psi_A(j)<E_A(j)\le2\psi_A(j)$. Put $y_j=\psi_A(j)$.
The recurrence and $y_{u-1}<y_u$ imply
\[
 y_{u+2}-2y_u>(4A^2-2A-3)y_u>A.
\]
Since $R\ge u+2$,
$E_A(R)-E_A(u)>A>X$, the last inequality following from
$a=Y(X+1)>X$. The right side of \eqref{eq:gammadiff} is positive,
proving $\sigma>0$.

Finally,
\[
 \kappa=u+\delta\Delta,\quad
 \mu=W_*+a\kappa+\rho H,\quad
 \mu^2-\Delta\kappa^2=1,\quad \gamma=\rho+\sigma.
\]
These are exactly the deleted input relations. Every retained source
row is unchanged, so the original positive input interface is restored.
\end{proof}

The converse is immediate from the original input theorem: at $x=1$
the marker is $W_*$, and setting $\gamma=\rho+\sigma$ and deleting
the three witnesses preserves every retained equation. Thus the
fixed-marker specialization has exactly the same accepted proposition
as the original fixed-input construction. This proves correctness of
$\mathcal S_{16}$ through the ordinary projection theorem.

The restoration argument is deliberately ordered before power recovery,
population counts, mask typing, or machine decoding. It therefore does
not presuppose the conclusion that the modified equations already
encode a computation.

\section{Correctness: eliminating the strong root}
\label{sec:elimproof}

Consider first the normalized sixteen-variable marker polynomial
\begin{equation}
 P_{16}^{\mathrm n}
 =\nf\nm\nw\nk\nt\,[Q(V^2-y^2)+y^2]\,[\Delta f^2-Q]-\Delta,
 \qquad Q=\Delta^2i^2c^4.
\label{eq:normalizedP}
\end{equation}
Its $74$-operation circuit is the fixed-marker reduction of complete84.
At a zero one first divides the equality by the positive common factor
$\Delta$, then applies the integer-unit and sign arguments of the
normalized kernel. It would be incorrect to infer that factors whose
product is $\Delta$ are individually units without that normalization.

The complete92 norm trick can now eliminate $f$. We include the full
local algebra because its sign issue matters.

\begin{lemma}[Integral reconstruction]
\label{lem:root}
Every positive zero of $P_{15}$ reconstructs a unique nonzero integer
$f=A_*/B_*$, satisfying $f^2=D_*$ and the normalized auxiliary
equations. This integer is in fact positive.
\end{lemma}

\begin{proof}
First, $D_*>c>0$. If $R\ge0$, then $b_*=c+RD_*>0$; if
$R\le-1$, then $b_*\le c-D_*<0$. Hence $B_*=2QcTb_*$ never
vanishes, even before the equations are imposed.

At a zero of \eqref{eq:P15}, both $\nf\nm\nw\nk\nt$ and
$N+1$ are integer units. Since $B_*$ is even,
\[
 N+1\equiv A_*^2+1\pmod4.
\]
The residue is $1$ or $2$, and cannot be $-1$. Thus $N=0$ and the
other product is $1$. The rational number $f=A_*/B_*$ satisfies
$f^2=D_*\in\Z$ and is therefore an integer: reducing the fraction
to lowest terms shows its denominator must be $1$.

From $f^2=D_*$ one gets $\Delta f^2-Q=\Delta$.
Moreover $V=u_*f-b_*$, and
\begin{align*}
 QV^2-(Q-1)y^2-1
 &=Q(D_*u_*^2+b_*^2)-(Q-1)y^2-1-2Qu_*b_*f\\
 &=A_*-B_*f=0.
\end{align*}
Thus the reconstructed signed $f$ gives a zero of
\eqref{eq:normalizedP}.

Suppose $f<0$. The simultaneous change $(f,T)\mapsto(-f,-T)$
preserves all expressions in the reconstructed normalized sixteen-variable
source, because $f,T$ occur there through $f^2$ and $Tf$.
Now $f>0$ and $T<0$. Apply the signed version of
Lemma~\ref{lem:rank}, and then Lemma~\ref{lem:restore}. This produces
a zero of the literal original complete84 source with all its witnesses
positive except possibly $T$. The repository's source-specific signed
quotient theorem proves that every such original zero on a valid compiler
slice has $T>0$ \cite{signed}. This is a contradiction. The signed
quotient theorem is applied only after the original input interface has
been restored, so no positive-domain theorem is applied prematurely.
\end{proof}

Conversely, a positive zero of the normalized sixteen-variable source
satisfies $f^2=D_*$ and $A_*=B_*f$. Deleting $f$ gives $N=0$ and
\eqref{eq:P15}. Therefore $P_{15}$ and the normalized marker source
have the same positive solvability. Together with
Lemma~\ref{lem:restore} and the fixed-program compiler this proves
the fifteen-variable part of Theorem~\ref{thm:main}.

The signed-quotient lemma is a genuine mathematical dependency. An
algebraic square-root reconstruction alone does not prove that its
reconstructed root is positive. Its pinned proof and independent review
are included in \code{vendor/}.

\section{A fully determined coefficient slice}
\label{sec:constants}

This section removes the common ambiguity ``choose coefficients
representing the desired recursively enumerable set.'' Every coefficient
here is fixed by a literal finite table and finite arithmetic. The
algorithm does not search for a halting computation or a ZFC refutation.
The original table has SHA-256
\begin{center}\small\ttfamily
fdbcef0825886035c3091c8537a69535bc96f9e265f0ae45863543d40ae70090.
\end{center}
It is included as \code{source/k215.tm}.

\subsection{From the literal K215 table to a normalized machine}

K215 has $215$ nonhalting states, two symbols, blank symbol $0$, initial
state $0$, and external halt state $215$. Its $430$ instructions are
preserved exactly. We construct a machine with $436$ states and six
symbols. Its ordered states are
\[
 \mathrm{start},\mathrm{enter},\mathrm{erase},\mathrm{return},
 \mathrm{dead},\mathrm{halt},p_0,\ldots,p_{214},m_0,\ldots,m_{214}.
\]
The ordered alphabet is $0,1,2,3,4,5$. At physical position $-n$,
$n>0$, symbols $0,1,4,5$ respectively store the two-track pairs
$(0,0),(0,1),(1,0),(1,1)$ for K215 positions $(n,-n)$.
At the origin, symbols $2,3$ store the origin bit $0,1$.

\begin{center}
\begin{tabular}{llll}
\toprule
State/read & New state & Write & Move\\
\midrule
start/1 & enter & 2 & stay\\
enter/2 & erase & 2 & left\\
erase/1 & erase & 0 & left\\
erase/0 & return & 0 & right\\
return/0 & return & 0 & right\\
return/2 & $p_0$ & 2 & stay\\
\bottomrule
\end{tabular}
\end{center}

Unspecified prefix transitions go to dead, preserving the symbol and
staying. Dead and halt are absorbing. A state $p_q$ or $m_q$ simulates
K215 instruction $q$ on the selected track. Positive-track movement is
reversed physically, negative-track movement is unchanged. At an origin
marker, either side state reads and updates the origin bit, moves left,
and chooses the positive or negative side from K215's direction. A
K215 target state $215$ becomes halt. All cases are literal in
\code{artifacts/folded_machine.json}.

On a positive unary word of length $\ell$ ending at the origin,
initialization takes exactly $2\ell+2$ steps and leaves a blank simulated
tape, state $p_0$, and head at the origin. The invariant maps a K215 head
position $h$ to physical position $-|h|$ and the selected side to the
sign of $h$; each simulated transition preserves it. Induction therefore
proves halting equivalence. For the chosen ordinary input $x=1$, the
compiler's unary length is $4$. The normalized machine accepts that
input exactly when K215 halts.

\subsection{The exact finite set of windows}

The compiler alphabet has $a_t=2628$ tiles, ordered as follows:
two boundary tiles; all pairs of a tape symbol with a head label in
\{$\mathrm{None}$, the $436$ states\}, in lexicographic order;
then the four initial-row phase tiles $L,I,Q,R$ with payloads
$(0,\mathrm{None}),(1,\mathrm{None}),(1,\mathrm{start}),(0,\mathrm{None})$.
The radius-one predicate is the exact function \code{local_valid}
in \code{source/compiler_k215.py}. It is the pinned helical tableau
predicate with stationary first-step markers \cite{helical}.

The predicate preserves vertical boundaries, synchronizes horizontal
boundaries, checks first-row phases, checks the three-neighbor machine
update with at most one arriving head, forbids crossing a vertical
boundary, and requires a last-row head to be halted. Its complete
boundary cases are supplied as executable finite logic, with no call
to a machine-acceptance oracle. This finite function is part of the
coefficient definition.

Enumerate valid crosses $(c,l,r,d,u)$ lexicographically. For each, enumerate
the four corners $(nw,ne,sw,se)$ lexicographically, yielding row-major
windows $(nw,d,ne,l,c,r,sw,u,se)$. Move these two valid windows to the
front, in the displayed order, omitting their later occurrences:
\begin{align*}
 w_0&=(1,1,1,2625,2625,2626,439,439,878),\\
 w_1&=(1,1,1,2624,2625,2625,2,439,439).
\end{align*}
There are exactly
\begin{align*}
 n_{\times}&=43\,126\,116\,188\,949\quad\text{valid crosses},\\
 n_w&=n_{\times}a_t^4
      =2\,057\,035\,525\,715\,746\,135\,339\,455\,744
      \quad\text{windows}.
\end{align*}
The count is obtained by factoring the finite sum over left and right
neighbors according to which head, if any, they send into the center.
For example, when the center sends no stationary head, the interior
count is
\[
 L_0R_0U_{s,0}+R_0\sum_q L_qU_{s,q}+L_0\sum_qR_qU_{s,q}.
\]
The script treats initialization and boundary cases separately. It
records each center's count. The factorization proof and the exact
conventions are documented in \code{proofs/COMPILER.md}.

\subsection{Layout and coefficient formulas}

Set $A_c=2^{92}$. This satisfies $A_c>n_w+1$ and $A_c\equiv1\pmod5$.
Let
\begin{align*}
 h_0&=9a_t+4, & m_0&=184+30a_t+8,\\
 p_c&=\left\lceil\frac{n_w+9a_t+4-m_0}{2}\right\rceil,&
 h_c&=h_0+p_c,& m&=m_0+2p_c,\\
 e_*&=n_w+12a_t,& M&=m+6a_t-3.
\end{align*}
There is exactly one ignored dummy position. Important numerical
values are
\begin{center}\small
\begin{tabular}{lr}
\toprule
$h_0$ & $23\,656$\\
$p_c$ & $1\,028\,517\,762\,857\,873\,067\,669\,700\,184$\\
$h_c$ & $1\,028\,517\,762\,857\,873\,067\,669\,723\,840$\\
$m$ & $2\,057\,035\,525\,715\,746\,135\,339\,479\,400$\\
$e_*$ & $2\,057\,035\,525\,715\,746\,135\,339\,487\,280$\\
$M$ & $2\,057\,035\,525\,715\,746\,135\,339\,495\,165$\\
\bottomrule
\end{tabular}
\end{center}
Write
\[
 \mathsf H=51M+3a_t+1,\quad
 \mathsf T_1=105M+4a_t+2,\quad
 \mathsf T_2=159M+4a_t+3,\quad g=186M+4a_t+4.
\]
Use the fixed $b,L_c,d,B$ from \eqref{eq:radices}, and put $r_c=2^b$.
Define the finite set
\[
 \mathcal E=\{0,\ldots,n_w-1\}\cup
 \{n_w+3a_t,\ldots,n_w+12a_t-1\}\cup
 \{e_*\}\cup\{M,3M,9M,27M\}.
\]
For $0\le r,c<3$ and $0\le s<a_t$ let
\[
 j(r,c,s)=1+(3r+c)a_t+s,\qquad
 e(r,c,s)=n_w+3a_t+(2-r)3a_t+(2-c)a_t+s.
\]
The coefficient attached to each exponent in $\mathcal E$ is
\begin{align*}
 c_i&=1+\sum_{r,c=0}^{2}A_c^{j(r,c,w_i[3r+c])}
          +\sum_{t=1}^{4}A_c^{9a_t+t} &&(0\le i<n_w),\\
 c_{e(r,c,s)}&=A_c^{j(r,c,s)},&
 c_{3^tM}&=A_c^{9a_t+1+t}\quad(0\le t<4),\\
 c_{e_*}&=0,&
 \mu_c&=A_c-2+\sum_{j=1}^{h_c}A_c^j.
\end{align*}
All exponent ranges in these sums are fixed finite ranges, not
Diophantine variables. Finally define
\begin{align}
 D_C={}&r_c^{3a_t}+r_c^{\mathsf H+a_t}+r_c^{8M}+r_c^{24M}
       +\sum_{e\in\mathcal E}c_e
           (r_c^{\mathsf T_1-e}+r_c^{\mathsf T_2-e}),\label{eq:DC}\\
 D_R={}&r_c^{\mathsf H},\label{eq:DR}\\
 \mc={}&B-1-\sum_{e\in\mathcal E\setminus\{1\}}r_c^e-2r_c^{e_*},
 \label{eq:MC}\\
 M_F^{\mathrm{native}}={}&\mu_c(r_c^{\mathsf T_1}+r_c^{\mathsf T_2})
       +\sum_{r=0}^{1}\sum_{c=0}^{2}\sum_{s=0}^{a_t-1}r_c^{e(r,c,s)}
       \notag\\
 &+\sum_{r=0}^{2}\sum_{c=0}^{1}\sum_{s=0}^{a_t-1}r_c^{\mathsf H+e(r,c,s)}
       +r_c^{9M}+r_c^{27M}+4.\label{eq:MF}
\end{align}
Together with \eqref{eq:ports}, these specify every coefficient port
of the displayed equations.

\subsection{Why this is an admissible compiler instance}

The complete75 compiler allows a power-of-two clause radix greater than
$n_w+1$, extra zero clauses, and a sufficiently large power-of-five
inner radix exponent. The above choices obey these rules. In particular,
$L_c>g+27M$, so all duplicated bands and anchors are separated.
The sum of center coefficients is bounded by
\[
 \sum_e c_e\le(14n_w+9a_t+4)A_c^{h_0}.
\]
The stronger modified75 no-carry margins follow by choosing $b$ larger
than binary exponent bounds for both $2\mu_c+4$ and
\[
 4(m+3)\bigl(2\sum_e c_e+6\bigr)+8.
\]
The two bounds computed in the certificate are
$94\,623\,634\,182\,924\,322\,225\,614\,593\,282$ and
$2\,176\,541$, respectively; $b=5^{42}$ exceeds both.
Since $A_c\equiv1$ and $r_c\equiv2\pmod5$, direct evaluation of the
finite coefficient recipe gives $2D_C-D_R\equiv2\pmod5$. Thus the
optional high-monomial correction is unnecessary.
Both $b$ and $L_c$ are powers of five, as required for the inherited
actual-index dummy-control argument.

The modifications $-2r_c^{e_*}$ and $+4$ in
\eqref{eq:MC}--\eqref{eq:MF} belong to the complete75 compiler and
must not be replaced by the older complete76 masks. The padded clause
radix $2^{92}$ and conservative $b=5^{42}$ increase coefficient size;
they do not change the local predicate or accepted computation.

Thus the inherited fixed-program theorem applies to this actual machine
and these actual fixed constants. The complete chain is
\[
\begin{aligned}
 \text{positive Diophantine solution}
 &\ \longleftrightarrow\ \text{normalized machine accepts }1^4\\
 &\ \longleftrightarrow\ \text{K215 halts}\\
 &\ \longleftrightarrow\ \exists p\,\ZRef(p).
\end{aligned}
\]
The first equivalence uses the proved marker restoration; the second
uses the explicit fold-and-erase simulation; the third uses the K215
calibration.

\section{Exact degree and operation certificates}
\label{sec:degrees}

Formal maximum-degree propagation through a circuit can overestimate
degrees after cancellation. Here the uniform degrees are proved
algebraically, and finite-field univariate evaluations serve only as
additional diagnostics.

Give every witness degree $1$ and every fixed numeral degree $0$. Set
\[
 Q_0=(B-1)J,\quad k_0=\eta+\zeta,\quad
 H_0=Q_0-F-2Z-\alpha,\quad
 T_2=w(Q_0-F-Z-\alpha)-tQ_0.
\]
The leading homogeneous forms of $\nf,\nm,\nk,\nt,\nw$ are
respectively
\[
 -k_0^2w^2s^4Q_0^{14},\quad
 8\gamma k_0w^2s^3Q_0^{11},\quad
 -hwsQ_0^4,\quad T_2,\quad H_0^2.
\]
Their degrees are $22,18,7,2,2$. For the main norm the degree-$22$
terms cancel exactly; writing $g_0=\gamma H$ gives the identity
\[
 (X+ac+g_0)^2-(a^2+H)c^2
 =X^2+2Xg_0+g_0^2+2acX+2acg_0-Hc^2,
\]
which exposes its degree-$18$ leading term.

In $\mathcal S_{16}$, the auxiliary residual has degree $28$ and leading
form
\[
 k_0^2w^2s^4Q_0^{14}T^2f^4.
\]
The other residuals have degrees $22,18,22,7,2,1$, in the corresponding
displayed order. Hence the sum of squares has exact degree $56$ and
leading form
\[
 k_0^4w^4s^8Q_0^{28}T^4f^8.
\]
The coefficient of $\eta^4w^4s^8J^{28}T^4f^8$ is
$(B-1)^{28}\ne0$.

For $P_{15}$, $Q,D_*,b_*,u_*$ have degrees $46,34,38,6$.
$A_*$ has degree $122$, while $B_*$ has degree $90$.
Therefore $N=A_*^2-D_*B_*^2$ has degree $244$: the second term has
only degree $214$. The outer product has degree $51$. The leading
form of $P_{15}$ is
\begin{equation}
 8h\gamma i^{12}k_0^{27}w^{21}s^{48}Q_0^{177}
 (Q_0-F)^4H_0^2T_2.
\label{eq:P15leader}
\end{equation}
In particular the coefficient of
\[
 J^{184}h\gamma i^{12}\eta^{27}w^{21}s^{48}t
\]
is $-8(B-1)^{184}$, uniformly nonzero. Its degree is $295$.

The JSON sources list every binary gate, and the independent checker
reconstructs every source row from the pinned parent packets. The full
$82$-gate and $64$-gate schedules are reproduced in the appendices.
Additional product variants have degrees $157,125,101$ and
operation counts $74,75,76$ respectively; their full circuits and
leading-form proof are included in \code{artifacts/} and \code{proofs/}.

\section{Nonnegative and unrestricted integer versions}
\label{sec:domains}

\subsection{Nonnegative variables without changing their number}

If $\N=\{0,1,2,\ldots\}$ is the preferred Diophantine domain, replace
every positive witness $z$ by $n+1$, with $n\in\N$. This preserves
the number of variables and the total degree, and preserves solvability
in both directions. In particular it gives a fifteen-nonnegative-variable
polynomial equivalent to refutability.

\subsection{Three unrestricted integers for each positive witness}

Define
\begin{equation}
 \operatorname{Pos}(a,b,c)=1+a^2+b^2+c(c+1).
\label{eq:Pos}
\end{equation}
Its image on $\Z^3$ is exactly $\Zp$. Positivity is immediate since
$c(c+1)\ge0$ for every integer $c$. Conversely, for $z>0$, Legendre's
three-square theorem applies to $4z-3$, a positive integer congruent
to $1$ or $5$ modulo $8$. Modulo $4$, a sum of three squares equal
to $1$ has exactly one odd square. Writing the two even roots as
$2a,2b$ and the odd root as $2c+1$ yields
$z=1+a^2+b^2+c(c+1)$ \cite{legendre}.

Substitution of independent copies of \eqref{eq:Pos} gives
\begin{center}
\begin{tabular}{lrrr}
\toprule
Form over $\Z$ & Integer unknowns & Equations & Degree\\
\midrule
$P_{15}(\operatorname{Pos},\ldots,\operatorname{Pos})$ & 45 & 1 & 590\\
$\mathcal S_{16}(\operatorname{Pos},\ldots,\operatorname{Pos})$ & 48 & 7 & 56\\
$H_{16}(\operatorname{Pos},\ldots,\operatorname{Pos})$ & 48 & 1 & 112\\
\bottomrule
\end{tabular}
\end{center}
The degree doubles exactly: restricting each triple to $(a,0,0)$ gives
the substitution $z\mapsto1+a^2$, which sends a nonzero leading form
to a nonzero form of twice its degree.

\subsection{A completely explicit quartic alternative}

Introduce one signed integer register for each of the $64$ gates of
the ordinary comparison circuit and keep the $16$ original inputs.
For each input $z$ add three signed integers and the quadratic equation
$z=\operatorname{Pos}(a,b,c)$. Each gate equation is linear or
quadratic; add the seven output comparisons. This gives
\[
 16+64+3\cdot16=128\quad\text{integer unknowns},\qquad
 64+7+16=87\quad\text{quadratic equations}.
\]
Their sum of squared residuals is a single quartic polynomial. The
degree is exactly four, since a fresh positivity coordinate's fourth
power occurs with coefficient $1$. Solvability is unchanged: evaluating
the gates extends a solution uniquely on registers, and the gate
equations force the intended evaluation in the converse direction.
This is an explicit degree-oriented alternative, not a claim to improve
classical existential quartic bounds. Complete equation trees and the
conversion generator are included in \code{artifacts/domain_variants.json}
and \code{source/domain_variants.py}.

\section{What is proved, checked, and still to formalize}
\label{sec:formal}

Three different kinds of evidence must be distinguished.

\textbf{Mathematical proof.} The new fixed-marker restoration and
strong-root argument are mathematical proofs. The historical compiler,
Pell rank, and signed-quotient proofs are identified precisely and
included as pinned source notes. The fold-and-erase machine reduction
is proved by a finite transition invariant. The Diophantine and simulation parts use finite syntax, natural-number
induction, elementary number theory, and finite computation coding.
The K215 calibration additionally uses set-model arguments, described
below. All of these mathematical arguments are available in ZFC.

\textbf{Inherited formal endpoint.} The repository supplies Lean and
Rocq theorems calibrating K215 halting to ZFC refutation. Its README
reports complete builds and states the proof-assistant axiom reports.
Those are inherited results, not builds rerun for this submission.
An external Lean or Rocq theorem is not by itself an encoded ZFC
derivation. Nor does the absence of user axioms establish a general
conservativity theorem for the assistant's foundations. The internal
ZFC claim concerns the particular syntactic, simulation, and set-model
calibration arguments, which can be carried out in ZFC; it does not invoke a
global soundness theorem for Lean or Rocq.


The actual calibration first compares the finite \code{C222} calculus
with ZF using first-order completeness, a seed model, a bisimulation
quotient, and a well-founded submodel. It then compares ZF and ZFC
refutations using the constructible submodel of a supplied set model
of ZF, together with first-order soundness and completeness. These are
ordinary ZFC arguments about set-sized domains, their subsets, quotients,
and recursively defined satisfaction relations. They do not require a
model of the entire proof-assistant universe hierarchy or an inaccessible
cardinal. In particular, the argument never asserts unconditionally that
ZFC has a model. The exact intermediate files and the set-domain
interpretation are described in \code{proofs/FORMAL_STATUS.md}.

The source presentation matches standard ZFC: it has the usual ZF axiom
schemas and Zermelo's selector form of Choice. Partial-functional
Replacement is equivalent to the usual total-functional presentation
by adding a default value and then separating the genuine range.
The selector form of Choice is equivalent to the choice-function form
by using disjoint tagged copies of a set-indexed family. Finite schema
annotations make numerical proof checking primitive recursive. Thus the
repository's syntactic consistency statement is provably equivalent in
ZFC to that for a standard effective first-order presentation.

\textbf{New executable audits.} The supplied standard-library Python
checks reconstruct the exact literal source transformations, verify
dependency closure and operation counts, verify source hashes and
coefficient layout inequalities, and perform independent algebra and
degree diagnostics. The finite simulation checks cover $5160$ local
folded transitions and seven initialization lengths
$1,2,3,4,7,19,64$. Both marker windows pass the local predicate. A
separate exhaustive $18$-tile fixture checks all $1\,889\,568$ crosses
and finds $121\,165$ valid crosses, agreeing with the factored count
at every center. This fixture corroborates the count implementation;
it is not a proof by testing of the unbounded compiler theorem.

No new Lean or Rocq kernel proof of the complete Diophantine theorem
is delivered here. In particular, the repository's mathematical
compiler notes remain a proof dependency. A later formalization should
encode the new marker lemma and ordinary/normalized interfaces, then
connect them to a formally verified machine-to-polynomial compiler.
The exact coefficient recipe and finite source circuits supplied here
make that task concrete and version-stable.

The statement is also meaningful in nonstandard models of ZFC. Internally
quantified proof codes and finite computations need not be standard
externally; the proof must relate them by the same arithmetic coding.
It is not enough to argue only about standard computations. The use
of the complete ZFC proof, including uniform induction, syntactic
simulation, and set-model calibration, supports the internal equivalence
without an appeal to the truth of consistency.

If ZFC is consistent, the displayed equations have no standard
solutions. By G\"odel's second incompleteness theorem, under the usual
recursive-presentation hypotheses ZFC then cannot prove their lack of
solutions. Mere consistency does not by itself justify the stronger
claim that ZFC also cannot prove that a solution exists; a suitable
soundness assumption would address that different assertion.

\section{Reproduction, integration, and comparison}
\label{sec:reproduce}

The archive is self-contained for its arithmetic and source checks and
uses only Python's standard library. From its root run
\begin{verbatim}
python3 source/verify.py
python3 source/render_circuits.py
latexmk -pdf -interaction=nonstopmode -halt-on-error report.tex
\end{verbatim}
The first command regenerates the compiler and circuit outputs in a
temporary directory, checks them against the delivered files, verifies
the pinned vendor hashes, reruns the independent source audit, and
checks the domain variants. It does not expand the giant fixed
integers, enumerate all $n_w$ windows, search for a ZFC contradiction,
or rerun the inherited assistant builds.

\begin{center}\small
\begin{tabular}{p{0.31\textwidth}p{0.61\textwidth}}
\toprule
Path & Contents\\
\midrule
\code{report.tex}, \code{report.pdf} & This report and its editable source.\\
\code{source/} & Literal K215 table, coefficient compiler, circuit transformation,
independent audit, domain conversion, and reproduction driver.\\
\code{artifacts/} & Complete machine table, finite coefficient recipe,
all polynomial circuits and comparison systems, integer-domain variants,
and generated LaTeX gate tables.\\
\code{proofs/} & Detailed new proof, independent mathematical review,
coefficient proof, domain proof, and provenance notes.\\
\code{vendor/} & Pinned upstream proof notes, original circuit packets,
endpoint sources, manifest, and inherited notices.\\
\code{checks/} & Reproducible validation receipts and delivery manifest.\\
\bottomrule
\end{tabular}
\end{center}

A natural repository destination is a new research subdirectory under
\code{Computability/HilbertTenthProblem/Papers/research-wip/}, with
the existing native-stream-queue files cited by the supplied manifest.
No existing repository files need to be overwritten. The archive's
own copies of upstream notes are frozen for review; their relative
cross-links may refer to the original repository layout.

The universal complete84 and complete92 sources use respectively
$18$ and $17$ positive witnesses. The new fixed-marker reduction
removes two witnesses in each case and reduces each circuit by ten
operations. The ordinary seven-equation form also removes unnecessary
product finalization and compares strong-equation sides directly.
The savings arise from the fixed input, so they must not be reported
as new bounds for a varying universal input.

Classical work supplies small universal variable/degree pairs
\cite{jones1982}, and dedicated work represents provability
Diophantinely \cite{carlmoroz}. This contribution's distinctive content
is its exact K215 specialization, concrete finite coefficient recipe,
proved removal of the input block, short shared circuits, and explicit
review boundary. Its greatest remaining cost is coefficient size:
$B=2^{5^{85}}$ alone has $5^{85}+1$ binary digits. A reduction in
machine or tableau encoding size could improve this cost without
changing the number of Diophantine unknowns.

\clearpage
\appendix
\section{Literal circuit for the fifteen-variable polynomial}

Each $r_j$ in this appendix is a polynomial abbreviation. It is not an
extra quantified variable. The table is generated directly from the
delivered \code{fixed_marker82} JSON packet. All six named coefficient
ports have the fixed values specified in Section~\ref{sec:constants}.

\begingroup\small
\begin{longtable}{r@{\quad}l@{\hspace{16mm}}r@{\quad}l}
\caption{The exact 82-operation circuit for $P_{15}$.}\\
\toprule Gate & Polynomial abbreviation & Gate & Polynomial abbreviation \\ \midrule
\endfirsthead
\toprule Gate & Abbreviation (continued) & Gate & Abbreviation (continued) \\ \midrule
\endhead
\bottomrule\endfoot
1 & $r_{1}=\tau \cdot \tau$ & 42 & $r_{42}=r_{2} \cdot r_{30}$ \\
2 & $r_{2}=(B-1) \cdot J$ & 43 & $r_{43}=r_{42} + r_{31}$ \\
3 & $r_{3}=r_{2} + 1$ & 44 & $r_{44}=r_{4} - 1$ \\
4 & $r_{4}=r_{3} \cdot r_{3}$ & 45 & $r_{45}=r_{43} \cdot r_{44}$ \\
5 & $r_{5}=r_{4} \cdot r_{3}$ & 46 & $r_{46}=r_{3} \cdot M_F$ \\
6 & $r_{6}=w \cdot r_{3}$ & 47 & $r_{47}=M_C + r_{46}$ \\
7 & $r_{7}=s \cdot r_{5}$ & 48 & $r_{48}=r_{47} \cdot J$ \\
8 & $r_{8}=r_{6} \cdot r_{7}$ & 49 & $r_{49}=r_{45} + r_{48}$ \\
9 & $r_{9}=\eta + \zeta$ & 50 & $r_{50}=r_{41} - r_{49}$ \\
10 & $r_{10}=r_{9} \cdot r_{7}$ & 51 & $r_{51}=K + w$ \\
11 & $r_{11}=r_{8} \cdot r_{10}$ & 52 & $r_{52}=r_{51} \cdot r_{33}$ \\
12 & $r_{12}=r_{11} + r_{9}$ & 53 & $r_{53}=r_{52} + r_{30}$ \\
13 & $r_{13}=r_{11} \cdot r_{12}$ & 54 & $r_{54}=t \cdot r_{2}$ \\
14 & $r_{14}=r_{1} - r_{13}$ & 55 & $r_{55}=r_{53} - r_{54}$ \\
15 & $r_{15}=r_{10} + \eta$ & 56 & $r_{56}=i \cdot r_{27}$ \\
16 & $r_{16}=r_{8} + r_{7}$ & 57 & $r_{57}=r_{56} \cdot r_{56}$ \\
17 & $r_{17}=r_{15} \cdot r_{16}$ & 58 & $r_{58}=i \cdot r_{26}$ \\
18 & $r_{18}=r_{6} + r_{17}$ & 59 & $r_{59}=r_{58} \cdot r_{56}$ \\
19 & $r_{19}=4 \cdot r_{16}$ & 60 & $r_{60}=r_{59} + 1$ \\
20 & $r_{20}=r_{19} + 3$ & 61 & $r_{61}=r_{49} \cdot r_{60}$ \\
21 & $r_{21}=\gamma \cdot r_{20}$ & 62 & $r_{62}=r_{15} + r_{61}$ \\
22 & $r_{22}=r_{18} + r_{21}$ & 63 & $r_{63}=r_{15} \cdot T$ \\
23 & $r_{23}=r_{22} \cdot r_{22}$ & 64 & $r_{64}=r_{63} \cdot r_{63}$ \\
24 & $r_{24}=r_{16} \cdot r_{16}$ & 65 & $r_{65}=r_{60} \cdot r_{64}$ \\
25 & $r_{25}=r_{24} + r_{20}$ & 66 & $r_{66}=r_{62} \cdot r_{62}$ \\
26 & $r_{26}=r_{15} \cdot r_{15}$ & 67 & $r_{67}=r_{38} \cdot r_{50}$ \\
27 & $r_{27}=r_{25} \cdot r_{26}$ & 68 & $r_{68}=r_{67} \cdot r_{55}$ \\
28 & $r_{28}=r_{23} - r_{27}$ & 69 & $r_{69}=r_{65} + r_{66}$ \\
29 & $r_{29}=r_{14} \cdot r_{28}$ & 70 & $r_{70}=r_{69} - r_{39}$ \\
30 & $r_{30}=r_{3} - F$ & 71 & $r_{71}=r_{57} \cdot r_{70}$ \\
31 & $r_{31}=r_{30} - Z$ & 72 & $r_{72}=r_{71} + r_{39}$ \\
32 & $r_{32}=r_{31} - \alpha$ & 73 & $r_{73}=r_{72} - 1$ \\
33 & $r_{33}=r_{32} - 2d$ & 74 & $r_{74}=r_{57} \cdot r_{65}$ \\
34 & $r_{34}=r_{33} - Z$ & 75 & $r_{75}=r_{57} \cdot r_{66}$ \\
35 & $r_{35}=r_{34} - W_*$ & 76 & $r_{76}=r_{74} \cdot r_{75}$ \\
36 & $r_{36}=r_{35} \cdot r_{35}$ & 77 & $r_{77}=4 \cdot r_{76}$ \\
37 & $r_{37}=r_{36} + 1$ & 78 & $r_{78}=r_{73} \cdot r_{73}$ \\
38 & $r_{38}=r_{29} \cdot r_{37}$ & 79 & $r_{79}=r_{78} - r_{77}$ \\
39 & $r_{39}=y \cdot y$ & 80 & $r_{80}=r_{79} + 1$ \\
40 & $r_{40}=h \cdot r_{8}$ & 81 & $r_{81}=r_{68} \cdot r_{80}$ \\
41 & $r_{41}=r_{9} - r_{40}$ & 82 & $r_{82}=r_{81} - 1$ \\
\end{longtable}
The output polynomial is $r_{82}$.

\endgroup

\clearpage
\section{Literal shared circuit for the seven-equation system}

The same convention applies: the registers are abbreviations unless
one deliberately uses the quadratic lift of Section~\ref{sec:domains}.
The order of output comparisons in the machine-readable packet differs
from the display in \eqref{eq:S16}, but the seven equations are identical.

\begingroup\small
\begin{longtable}{r@{\quad}l@{\hspace{16mm}}r@{\quad}l}
\caption{The exact 64-operation circuit for the seven comparisons.}\\
\toprule Gate & Polynomial abbreviation & Gate & Polynomial abbreviation \\ \midrule
\endfirsthead
\toprule Gate & Abbreviation (continued) & Gate & Abbreviation (continued) \\ \midrule
\endhead
\bottomrule\endfoot
1 & $r_{1}=\tau \cdot \tau$ & 33 & $r_{33}=h \cdot r_{25}$ \\
2 & $r_{2}=(B-1) \cdot J$ & 34 & $r_{34}=r_{28} + r_{24}$ \\
3 & $r_{3}=\eta + \zeta$ & 35 & $r_{35}=r_{6} \cdot r_{27}$ \\
4 & $r_{4}=f \cdot f$ & 36 & $r_{36}=r_{29} + r_{3}$ \\
5 & $r_{5}=y \cdot y$ & 37 & $r_{37}=r_{30} \cdot r_{31}$ \\
6 & $r_{6}=K + w$ & 38 & $r_{38}=4 \cdot r_{31}$ \\
7 & $r_{7}=T \cdot f$ & 39 & $r_{39}=r_{31} \cdot r_{31}$ \\
8 & $r_{8}=r_{2} + 1$ & 40 & $r_{40}=r_{30} \cdot r_{30}$ \\
9 & $r_{9}=r_{4} - 1$ & 41 & $r_{41}=r_{3} - r_{33}$ \\
10 & $r_{10}=t \cdot r_{2}$ & 42 & $r_{42}=r_{35} + r_{14}$ \\
11 & $r_{11}=r_{7} - 1$ & 43 & $r_{43}=r_{30} \cdot r_{11}$ \\
12 & $r_{12}=r_{8} \cdot r_{8}$ & 44 & $r_{44}=r_{34} \cdot r_{4}$ \\
13 & $r_{13}=w \cdot r_{8}$ & 45 & $r_{45}=r_{29} \cdot r_{36}$ \\
14 & $r_{14}=r_{8} - F$ & 46 & $r_{46}=r_{13} + r_{37}$ \\
15 & $r_{15}=r_{8} \cdot M_F$ & 47 & $r_{47}=r_{38} + 3$ \\
16 & $r_{16}=r_{12} \cdot r_{8}$ & 48 & $r_{48}=r_{41} - r_{34}$ \\
17 & $r_{17}=r_{14} - Z$ & 49 & $r_{49}=r_{42} - r_{10}$ \\
18 & $r_{18}=r_{2} \cdot r_{14}$ & 50 & $r_{50}=i \cdot r_{40}$ \\
19 & $r_{19}=r_{12} - 1$ & 51 & $r_{51}=r_{43} - r_{44}$ \\
20 & $r_{20}=M_C + r_{15}$ & 52 & $r_{52}=r_{1} - r_{45}$ \\
21 & $r_{21}=s \cdot r_{16}$ & 53 & $r_{53}=\gamma \cdot r_{47}$ \\
22 & $r_{22}=r_{17} - \alpha$ & 54 & $r_{54}=r_{39} + r_{47}$ \\
23 & $r_{23}=r_{18} + r_{17}$ & 55 & $r_{55}=r_{51} \cdot r_{51}$ \\
24 & $r_{24}=r_{20} \cdot J$ & 56 & $r_{56}=r_{50} \cdot r_{50}$ \\
25 & $r_{25}=r_{13} \cdot r_{21}$ & 57 & $r_{57}=r_{46} + r_{53}$ \\
26 & $r_{26}=r_{3} \cdot r_{21}$ & 58 & $r_{58}=r_{54} \cdot r_{40}$ \\
27 & $r_{27}=r_{22} - 2d$ & 59 & $r_{59}=r_{54} \cdot r_{9}$ \\
28 & $r_{28}=r_{23} \cdot r_{19}$ & 60 & $r_{60}=r_{55} - r_{5}$ \\
29 & $r_{29}=r_{25} \cdot r_{26}$ & 61 & $r_{61}=r_{57} \cdot r_{57}$ \\
30 & $r_{30}=r_{26} + \eta$ & 62 & $r_{62}=r_{59} \cdot r_{60}$ \\
31 & $r_{31}=r_{25} + r_{21}$ & 63 & $r_{63}=r_{61} - r_{58}$ \\
32 & $r_{32}=r_{27} - Z$ & 64 & $r_{64}=r_{62} + r_{5}$ \\
\end{longtable}
The seven comparisons are $r_{52}=1,\ r_{63}=1,\ r_{48}=1,\ r_{49}=1,\ r_{64}=1,\ r_{56}=r_{59},\ r_{32}=W_*$.

\endgroup

\clearpage
\begingroup\small
\begin{thebibliography}{99}
\addcontentsline{toc}{section}{References and pinned proof dependencies}
\bibitem{repo}
Vladimir Reshetnikov, \emph{ProveIt}, source snapshot
\commit. \repo.
The full repository paths and SHA-256 values of every included upstream
file appear in \code{vendor/manifest.json}.

\bibitem{k215}
ProveIt, \emph{K215 verification README} and
\emph{KEndpoint215.lean}, at the pinned snapshot.
\href{https://github.com/VladimirReshetnikov/ProveIt/tree/cdd582bec13f908df3711fa9e05d1ddac368a5db/Computability/BusyBeaver/Research/ZFC/k225-formal/verification/k215}{Pinned K215 verification directory}.
The source endpoint is
\code{SetTheory.BusyBeaver.ZFC.KEndpoint215.blank_halts_iff_zfc_refutation}.

\bibitem{compiler}
ProveIt, \code{FIXED_RAW_UNIVERSAL_75_PROOF.md},
\code{FIXED_RAW_UNIVERSAL_76_PROOF.md},
and \code{complete75_half_binomial_compiler.md}, pinned snapshot.
The associated layout, no-carry, and actual-index arguments are included
with their immediate proof dependencies in \code{vendor/}.

\bibitem{helical}
ProveIt, \code{FIXED_RAW_UNIVERSAL_78_PROOF.md},
\code{FIXED_RAW_UNIVERSAL_81_PROOF.md}, and
\code{explore_helical_unary_tableau.py}, pinned snapshot.
These define the native window compiler and stationary input markers.

\bibitem{pell}
ProveIt, \code{review_complete85_auxiliary_bezout_math.md},
\code{review_complete86_ordinary_auxiliary_math.md},
\code{PELL_RELAXED_AUXILIARY_PROOF.md}, and
\code{complete75_gamma_dominance_elimination102.md}, pinned snapshot.

\bibitem{projections}
ProveIt, the complete84 scaled-strong, complete85 reduced-auxiliary,
complete86 ordinary-auxiliary, and complete92 norm-unit-guard source
packets and accompanying mathematical notes, pinned snapshot.
Exact input JSON bytes are included in \code{vendor/}.

\bibitem{signed}
ProveIt, \code{complete84_signed_quotient_absorption.md},
\code{complete84_signed_quotient_soundness_scout.md}, and
\code{review_complete84_signed_quotient_soundness.md}, pinned snapshot.

\bibitem{jones1982}
James P. Jones, ``Universal Diophantine equation,''
\emph{Journal of Symbolic Logic} \textbf{47} (1982), 549--571.
\href{https://doi.org/10.2307/2273588}{doi:10.2307/2273588}.

\bibitem{carlmoroz}
M. Carl and B. Z. Moroz, ``On a Diophantine Representation
of the Predicate of Provability,'' \emph{Journal of Mathematical
Sciences} \textbf{199} (2014), 36--52.
\href{https://doi.org/10.1007/s10958-014-1830-2}{doi:10.1007/s10958-014-1830-2}.

\bibitem{legendre}
A.-M. Legendre, \emph{Essai sur la th\'eorie des nombres}, Paris,
1798. Three-square theorem: a positive integer is a sum of three
integer squares if and only if it is not of the form $4^a(8b+7)$.
Only this classical theorem is used for the optional unrestricted-integer
conversion; the principal positive-variable constructions do not need it.
\end{thebibliography}
\endgroup

\end{document}
